\chapter{Barisan: Perkenalan Pertama}\label{ch:g11:seq}

Sebuah \omterm{def:g11:seq:sequence}{barisan} adalah daftar bilangan yang dihasilkan oleh sebuah aturan: saldo
tabungan yang berturut-turut, ukuran sebuah populasi dari tahun ke tahun. Bab
ini menelaah dua keluarga yang menguasai penerapannya ---
\omterm{def:g11:seq:sequence}{barisan} \emph{\omterm{def:g11:seq:arithmetic}{aritmetika}}, yang tumbuh dengan langkah yang sama, dan
\omterm{def:g11:seq:sequence}{barisan} \emph{\omterm{def:g11:seq:geometric}{geometri}}, yang tumbuh dengan \omterm{def:g11:seq:geometric}{rasio} yang sama. Teori limit yang
cermat dikembangkan pada \cref{ch:g12:seq}.

\section{Mendefinisikan sebuah barisan}

\begin{definition}[Barisan]\label{def:g11:seq:sequence}
Sebuah \emph{barisan}\index{barisan} $(u_n)$ memasangkan setiap \omterm{def:g10:numbers:sets}{bilangan bulat}
$n \geq 0$ (atau $n \geq 1$) dengan sebuah \omterm{def:g10:numbers:sets}{bilangan real} $u_n$, yang disebut
\emph{suku ke-$n$}. Sebuah barisan dapat diberikan
\begin{itemize}
  \item secara \emph{eksplisit}, lewat rumus $u_n$ dalam $n$: misalnya
        $u_n = n^2 + 1$;
  \item secara \emph{rekursif}, lewat suku pertamanya dan aturan untuk beralih
        dari setiap suku ke suku berikutnya: misalnya $u_0 = 3$ dan $u_{n+1} = 2u_n - 1$.
\end{itemize}
\end{definition}

\begin{example}
Untuk $u_n = n^2 + 1$: $u_0 = 1$, $u_1 = 2$, $u_2 = 5$, dan $u_{10} = 101$
secara langsung. Untuk $u_0 = 3$, $u_{n+1} = 2u_n - 1$: $u_1 = 5$, $u_2 = 9$,
$u_3 = 17$ --- setiap sukunya memerlukan suku sebelumnya; mencapai $u_{10}$
menuntut sepuluh langkah (atau sebuah rumus umum, lihat \cref{exo:g11:seq:11}).
\end{example}

\section{Barisan aritmetika}

\begin{definition}[Barisan aritmetika]\label{def:g11:seq:arithmetic}
Sebuah \omterm{def:g11:seq:sequence}{barisan} disebut \emph{aritmetika}\index{barisan aritmetika} dengan
\emph{beda}\index{beda} $d$ bila setiap sukunya diperoleh dari suku
sebelumnya dengan menambahkan $d$:
\[
u_{n+1} = u_n + d \quad \text{untuk semua } n.
\]
Setara dengan itu: selisih $u_{n+1} - u_n$ tetap, yaitu sama dengan $d$.
\end{definition}

\begin{theorem}[Suku umum]\label{thm:g11:seq:arithgeneral}
Jika $(u_n)$ \omterm{def:g11:seq:arithmetic}{aritmetika} dengan suku pertama $u_0$ dan \omterm{def:g11:seq:arithmetic}{beda} $d$,
maka
\[
u_n = u_0 + n\,d \quad \text{untuk semua } n \geq 0,
\qquad\text{dan secara lebih umum } u_n = u_p + (n - p)\,d .
\]
\end{theorem}

\begin{proof}
Untuk beralih dari $u_0$ ke $u_n$, aturan ``tambahkan $d$'' diterapkan $n$
kali: satu langkah memberi $u_1 = u_0 + d$, dua langkah memberi
$u_2 = u_0 + 2d$, dan setelah $n$ langkah setiap penerapannya telah
menyumbang satu $d$, sehingga $u_n = u_0 + nd$. (Kata ``dan seterusnya'' ini
dipertegas lewat induksi pada \cref{ch:g12:seq}.) Rumus umumnya menyusul
dengan membilang $n - p$ langkah dari $u_p$ ke $u_n$.
\end{proof}

\begin{theorem}[Jumlah bilangan bulat berurutan]\label{thm:g11:seq:intsum}
Untuk setiap \omterm{def:g10:numbers:sets}{bilangan bulat} $n \geq 1$:
\[
1 + 2 + \dots + n = \frac{n(n+1)}{2}.
\]
Secara lebih umum, jumlah suku berurutan sebuah \omterm{def:g11:seq:arithmetic}{barisan aritmetika} sama
dengan
\[
(\text{banyaknya suku}) \times
\frac{\text{suku pertama} + \text{suku terakhir}}{2}.
\]
\end{theorem}

\begin{proof}
Tulislah jumlahnya $S$ dua kali, yang kedua dalam urutan terbalik, lalu
jumlahkan kolom demi kolom:
\[
\begin{array}{ccccccccc}
S & = & 1 & + & 2 & + & \dots & + & n\\
S & = & n & + & (n-1) & + & \dots & + & 1\\
\hline
2S & = & (n+1) & + & (n+1) & + & \dots & + & (n+1)
\end{array}
\]
Ada $n$ kolom, masing-masing berjumlah $n + 1$, sehingga $2S = n(n+1)$. Untuk
\omterm{def:g11:seq:arithmetic}{barisan aritmetika} umum, pemasangan yang sama juga berhasil: pertama $+$
terakhir $=$ kedua $+$ kedua dari belakang $= \dots$, sebab bergerak satu
langkah maju di ujung kiri ($+d$) diimbangi oleh satu langkah mundur di ujung
kanan ($-d$).
\end{proof}

\begin{example}
$1 + 2 + \dots + 100 = \frac{100 \times 101}{2} = 5050$. Jumlah
bilangan ganjil $1 + 3 + \dots + 99$ ($50$ suku) adalah
$50 \times \frac{1 + 99}{2} = 2500$.
\end{example}

\section{Barisan geometri}

\begin{definition}[Barisan geometri]\label{def:g11:seq:geometric}
Sebuah \omterm{def:g11:seq:sequence}{barisan} disebut \emph{geometri}\index{barisan geometri} dengan
\emph{rasio}\index{rasio} $q \neq 0$ bila setiap sukunya diperoleh dari suku
sebelumnya dengan mengalikan $q$:
\[
u_{n+1} = q\,u_n \quad \text{untuk semua } n.
\]
Setara dengan itu, bila tidak ada sukunya yang lenyap: rasio
$\frac{u_{n+1}}{u_n}$ tetap, yaitu sama dengan $q$.
\end{definition}

\begin{theorem}[Suku umum]\label{thm:g11:seq:geomgeneral}
Jika $(u_n)$ \omterm{def:g11:seq:geometric}{geometri} dengan suku pertama $u_0$ dan \omterm{def:g11:seq:geometric}{rasio} $q$, maka
\[
u_n = u_0\, q^n \quad \text{untuk semua } n \geq 0,
\qquad\text{dan secara lebih umum } u_n = u_p\, q^{\,n-p} .
\]
\end{theorem}

\begin{proof}
Pembilangan langkah yang sama seperti pada \cref{thm:g11:seq:arithgeneral}:
dari $u_0$ ke $u_n$, aturan ``kalikan dengan $q$'' diterapkan $n$ kali, dan
menyumbang faktor $q^n$.
\end{proof}

\begin{theorem}[Jumlah geometri]\label{thm:g11:seq:geomsum}
Untuk setiap \omterm{def:g10:numbers:sets}{bilangan real} $q \neq 1$ dan \omterm{def:g10:numbers:sets}{bilangan bulat} $n \geq 0$:
\[
1 + q + q^2 + \dots + q^n = \frac{1 - q^{\,n+1}}{1 - q}.
\]
\end{theorem}

\begin{proof}
Misalkan $S = 1 + q + \dots + q^n$. Kalikan dengan $q$:
$qS = q + q^2 + \dots + q^{n+1}$. Lalu kurangkan:
\[
S - qS = \bigl(1 + q + \dots + q^n\bigr)
- \bigl(q + q^2 + \dots + q^{n+1}\bigr) = 1 - q^{\,n+1},
\]
sebab setiap suku antaranya muncul sekali pada masing-masing jumlah lalu
saling menghapus. Jadi $(1 - q)S = 1 - q^{\,n+1}$, dan membaginya dengan
$1 - q \neq 0$ memberi rumusnya.
\end{proof}

\begin{example}\label{ex:g11:seq:chessboard}
$1 + 2 + 4 + \dots + 2^{10} = \frac{1 - 2^{11}}{1 - 2} = 2^{11} - 1 =
2047$: melipatduakan butir beras pada petak papan catur mengalahkan
lumbung mana pun, jauh sebelum petak ke-$64$, yang di sana jumlahnya
$2^{64} - 1 \approx 1.8 \times 10^{19}$.
\end{example}

\begin{omfigure}
\begin{tikzpicture}
\begin{axis}[omaxis,
  width=10cm, height=5.6cm,
  xlabel={$n$}, ylabel={$u_n$},
  xmin=-0.5, xmax=10.8, ymin=0, ymax=13,
  xtick={5,10}, ytick={2,6}]
  \addplot[omDef, only marks, mark=*, mark size=1.5pt,
    samples at={0,...,10}] {2 + 0.9*x};
  \addplot[omThm, only marks, mark=*, mark size=1.5pt,
    samples at={0,...,10}] {2*1.2^x};
  \node[omDef, font=\small, anchor=east] at (axis cs:9.7,10.9)
    {aritmetika};
  \node[omThm, font=\small, anchor=west] at (axis cs:8.6,7.1)
    {geometri};
\end{axis}
\end{tikzpicture}

{\small Langkah yang sama melawan \omterm{def:g11:seq:geometric}{rasio} yang sama: \omterm{def:g11:seq:arithmetic}{barisan aritmetika}
($u_{n+1} = u_n + 0.9$, biru) mengikuti sebuah garis, \omterm{def:g11:seq:geometric}{barisan geometri}
($u_{n+1} = 1.2\,u_n$, merah) mengikuti kurva eksponen yang akhirnya
melampauinya.}
\end{omfigure}

\begin{method}[Mengenali jenis sebuah barisan]\label{met:g11:seq:recognize}
Hitunglah $u_{n+1} - u_n$ lalu sederhanakan. Jika hasilnya sebuah konstanta
$d$, \omterm{def:g11:seq:sequence}{barisannya} \omterm{def:g11:seq:arithmetic}{aritmetika}. Bila bukan, hitunglah $\frac{u_{n+1}}{u_n}$
(sukunya taknol) lalu sederhanakan: sebuah konstanta $q$ berarti \omterm{def:g11:seq:geometric}{geometri}.
Jika keduanya tidak tetap, \omterm{def:g11:seq:sequence}{barisannya} bukan salah satu jenis itu --- jangan
pernah menyimpulkan hanya dari beberapa suku pertamanya.
\end{method}

\begin{example}
Untuk $u_n = 3 \times 5^n$:
$\frac{u_{n+1}}{u_n} = \frac{3 \times 5^{n+1}}{3 \times 5^n} = 5$ untuk semua
$n$: \omterm{def:g11:seq:geometric}{geometri} dengan \omterm{def:g11:seq:geometric}{rasio} $5$. Untuk $u_n = n^2$:
$u_1 - u_0 = 1$ tetapi $u_2 - u_1 = 3$, dan
$\frac{u_1}{u_0}$ bahkan tidak terdefinisi --- jadi bukan \omterm{def:g11:seq:arithmetic}{aritmetika} maupun
\omterm{def:g11:seq:geometric}{geometri}.
\end{example}

\section{Kemonotonan}

\begin{definition}[Barisan monoton]\label{def:g11:seq:monotone}
Sebuah \omterm{def:g11:seq:sequence}{barisan} $(u_n)$ disebut \emph{naik} bila $u_{n+1} \geq u_n$ untuk
semua $n$, dan \emph{turun} bila $u_{n+1} \leq u_n$ untuk semua $n$.
\end{definition}

\begin{method}[Menelaah kemonotonan]\label{met:g11:seq:monotonicity}
Telaahlah tanda $u_{n+1} - u_n$. Untuk \omterm{def:g11:seq:sequence}{barisan} yang sukunya positif, kita
boleh membandingkan $\frac{u_{n+1}}{u_n}$ dengan $1$ sebagai gantinya.
\end{method}

\begin{example}
Sebuah \omterm{def:g11:seq:arithmetic}{barisan aritmetika} naik bila $d \geq 0$
($u_{n+1} - u_n = d$), dan turun bila $d \leq 0$. Sebuah \omterm{def:g11:seq:geometric}{barisan geometri}
dengan $u_0 > 0$ dan $q > 1$ bersifat naik:
$u_{n+1} - u_n = u_0 q^n (q - 1) > 0$; dengan $u_0 > 0$ dan $0 < q < 1$
\omterm{def:g11:seq:sequence}{barisan} itu turun.
\end{example}

\section{Perilaku jangka panjang, secara takresmi}

Apa yang terjadi pada $u_n$ ketika $n$ menjadi sangat besar? Untuk \omterm{def:g11:seq:arithmetic}{barisan
aritmetika} dengan $d > 0$, suku $u_0 + nd$ akhirnya melampaui bilangan tetap
mana pun. Untuk \omterm{def:g11:seq:geometric}{barisan geometri} dengan $0 < q < 1$, suku
$u_0 q^n$ mengerut menuju $0$: mengalikan berulang kali dengan $0.9$, misalnya,
mengikis nilai awal berapa pun. Dan untuk $q > 1$ sukunya meledak, seperti pada
\cref{ex:g11:seq:chessboard}.

\begin{remark}
Pernyataan ini dapat dibuat benar-benar tepat --- ``sukunya akhirnya
tetap berada dalam jarak berapa pun dari $0$'' --- lalu dibuktikan. Itulah teori
\emph{limit}, tema pembuka \cref{ch:g12:seq}.
\end{remark}

\section{Latihan}

\begin{exercise}[$\star$]\label{exo:g11:seq:1}
Untuk setiap \omterm{def:g11:seq:sequence}{barisan} berikut, hitunglah $u_1$, $u_2$, $u_3$:
\[
u_n = \frac{n}{n+1}; \qquad
u_0 = 5,\ u_{n+1} = 3u_n - 2; \qquad
u_n = (-1)^n\,n .
\]
\end{exercise}

\begin{exercise}[$\star$]\label{exo:g11:seq:2}
$(u_n)$ \omterm{def:g11:seq:arithmetic}{aritmetika} dengan $u_0 = 7$ dan $d = -3$. Hitunglah $u_{10}$ dan
$u_{25}$. $(v_n)$ \omterm{def:g11:seq:arithmetic}{aritmetika} dengan $v_3 = 11$ dan $v_8 = 26$. Carilah
\omterm{def:g11:seq:arithmetic}{bedanya} dan $v_0$.
\end{exercise}

\begin{exercise}[$\star$]\label{exo:g11:seq:3}
$(u_n)$ \omterm{def:g11:seq:geometric}{geometri} dengan $u_0 = 5$ dan $q = 2$. Hitunglah $u_8$. $(v_n)$
\omterm{def:g11:seq:geometric}{geometri} dengan suku positif, $v_2 = 12$ dan $v_4 = 48$. Carilah \omterm{def:g11:seq:geometric}{rasionya}
dan $v_0$.
\end{exercise}

\begin{exercise}[$\star$]\label{exo:g11:seq:4}
Hitunglah
\[
1 + 2 + 3 + \dots + 500, \qquad
4 + 7 + 10 + \dots + 61, \qquad
1 + \frac12 + \frac14 + \dots + \frac{1}{2^{10}} .
\]
\end{exercise}

\begin{exercise}[$\star$]\label{exo:g11:seq:5}
Tentukan apakah setiap \omterm{def:g11:seq:sequence}{barisan} berikut \omterm{def:g11:seq:arithmetic}{aritmetika}, \omterm{def:g11:seq:geometric}{geometri}, atau bukan keduanya:
\[
u_n = 4n - 1; \qquad
v_n = \frac{2^n}{3^{n+1}}; \qquad
w_n = n^2 + n .
\]
\end{exercise}

\begin{exercise}[$\star\star$]\label{exo:g11:seq:6}
Sebuah gedung pertunjukan mempunyai $20$ baris: $16$ kursi pada baris pertama,
dan setiap baris mempunyai $2$ kursi lebih banyak daripada baris sebelumnya.
Berapa kursi pada baris terakhir? Berapa kursi di seluruh gedung?
\end{exercise}

\begin{exercise}[$\star\star$]\label{exo:g11:seq:7}
Sebuah populasi bakteri berlipat dua setiap jam; pada tengah hari ada $500$
bakteri. Berapa banyak pada pukul 20.00? Setelah berapa jam penuh populasinya
pertama kali melampaui satu juta? (Selesaikan dengan mencoba pangkat
$2$ berturut-turut.)
\end{exercise}

\begin{exercise}[$\star\star$]\label{exo:g11:seq:8}
Setiap bulan, seorang penabung menyetorkan $100$ euro ke sebuah rekening yang
memberi bunga $0.2\%$ per bulan atas saldo yang ada (bunganya dibukukan
tepat sebelum setoran). Misalkan $c_n$ saldo tepat sesudah setoran ke-$n$,
sehingga $c_1 = 100$ dan $c_{n+1} = 1.002\,c_n + 100$. Hitunglah $c_2$
dan $c_3$, lalu jelaskan mengapa $(c_n)$ bukan \omterm{def:g11:seq:arithmetic}{aritmetika} maupun \omterm{def:g11:seq:geometric}{geometri}.
\end{exercise}

\begin{exercise}[$\star\star$]\label{exo:g11:seq:9}
Telaahlah kemonotonan \omterm{def:g11:seq:sequence}{barisan}
\[
u_n = n^2 - 8n \ (n \geq 0), \qquad
v_n = \frac{3^n}{n!}\ (n \geq 1),
\]
dengan $n! = 1 \times 2 \times \dots \times n$. (Untuk $(v_n)$, bandingkan
$\frac{v_{n+1}}{v_n}$ dengan $1$.)
\end{exercise}

\begin{exercise}[$\star\star$]\label{exo:g11:seq:10}
Jumlah $n$ suku pertama sebuah \omterm{def:g11:seq:arithmetic}{barisan aritmetika} dengan $u_0 = 3$
dan $d = 4$ sama dengan $903$. Carilah $n$. (Susunlah sebuah \omterm{def:g10:algebra:equation}{persamaan} kuadrat
dalam $n$ lalu pakai \cref{ch:g11:quad}.)
\end{exercise}

\begin{exercise}[$\star\star\star$]\label{exo:g11:seq:11}
Misalkan $u_0 = 3$ dan $u_{n+1} = 2u_n - 1$.
\begin{enumerate}
  \item Hitunglah $u_1, u_2, u_3$ lalu dugalah sebuah rumus untuk $u_n$.
  \item Misalkan $v_n = u_n - 1$. Tunjukkan bahwa $(v_n)$ \omterm{def:g11:seq:geometric}{geometri}, lalu
        berikan \omterm{def:g11:seq:geometric}{rasionya} dan suku pertamanya.
  \item Simpulkan rumus eksplisit untuk $u_n$ lalu periksalah dugaanmu.
\end{enumerate}
\end{exercise}

\section{Soal: Menara Brahma dan kelinci Fibonacci}

\begin{problem}[{Soal akhir pekan --- dua rekurensi legendaris:
menara yang mengakhiri dunia, barisan yang tumbuh bak emas,
dan kiat barisan bantu yang menjinakkan pinjaman}]
\label{pb:g11:seq:1}
Dua \omterm{def:g11:seq:sequence}{barisan} menguasai cerita rakyat matematika. Yang satu membilang
langkah pemindahan Menara Brahma --- enam puluh empat cakram emas yang,
menurut legenda, pemindahannya akan mengakhiri dunia. Yang lain membilang
kelinci Fibonacci dan menyembunyikan \omterm{def:g11:seq:geometric}{rasio} emas. Tidak satu pun
\omterm{def:g11:seq:arithmetic}{aritmetika}, tidak satu pun \omterm{def:g11:seq:geometric}{geometri} --- dan keduanya menyerah kepada
senjata bab ini: rekurensi, jumlah \omterm{def:g11:seq:geometric}{geometri}
(\cref{thm:g11:seq:geomsum}), dan kiat \omterm{def:g11:seq:sequence}{barisan} bantu pada
\cref{exo:g11:seq:11}, yang juga menghitung cicilan rumahmu.

\medskip
\noindent\textbf{Bagian I --- Menara Brahma.}
Teka-tekinya: $n$ cakram yang ukurannya mengecil ditumpuk pada tiang A;
pindahkan seluruh tumpukan ke tiang C, satu cakram setiap kali, tanpa pernah
meletakkan cakram yang lebih besar di atas yang lebih kecil (tiang B boleh
membantu). Misalkan $h_n$ banyaknya langkah minimal.
\begin{enumerate}
  \item Mainkan (dengan uang logam) lalu catat $h_1$, $h_2$, $h_3$.
  \item Jelaskan siasat di balik rekurensi
        $h_{n+1} = 2h_n + 1$: apa yang harus terjadi sebelum dan sesudah
        cakram terbesar berpindah?
  \item Selesaikan rekurensinya dengan kiat
        \cref{exo:g11:seq:11}: ambil $v_n = h_n + 1$, tunjukkan
        $(v_n)$ \omterm{def:g11:seq:geometric}{geometri}, lalu simpulkan $h_n = 2^n - 1$.
  \item Menara dalam legenda itu mempunyai $64$ cakram, dan para biarawannya
        memindahkan satu cakram per detik. Dengan memakai $2^{10} = 1024 \approx 10^3$,
        taksirlah waktu pemindahannya dalam tahun (setahun kira-kira
        $3 \times 10^7$ detik; bandingkan dengan
        \cref{ex:g11:seq:chessboard}, raksasa yang sama dalam cerita
        yang lain). Perlukah kita cemas?
  \item Mengapa tidak ada siasat yang dapat mengalahkan $2^n - 1$ langkah?
        Berargumenlah bahwa \emph{setiap} penyelesaian menuruti
        $h_{n+1} \geq 2 h_n + 1$: apa yang harus benar tentang $n$ cakram
        teratas tepat sebelum, dan tepat sesudah, cakram terbawah
        berpindah?
\end{enumerate}

\medskip
\noindent\textbf{Bagian II --- Fibonacci.}
Definisikan $F_1 = F_2 = 1$ dan $F_{n+2} = F_{n+1} + F_n$ (setiap sukunya
jumlah dua suku sebelumnya --- aturan pencacahan irama di jilid
sebelumnya, kini dengan nama Eropanya).
\begin{enumerate}[resume]
  \item Daftarkan $F_1$ sampai $F_{12}$.
  \item Tunjukkan bahwa $(F_n)$ bukan \omterm{def:g11:seq:arithmetic}{aritmetika} maupun \omterm{def:g11:seq:geometric}{geometri},
        tetapi tegas naik mulai dari $n = 2$
        (\cref{met:g11:seq:monotonicity} dan rekurensinya).
  \item Buktikan identitas jumlahnya
        \[
        F_1 + F_2 + \dots + F_n = F_{n+2} - 1
        \]
        secara teleskopik: tulislah setiap $F_k$ sebagai
        $F_{k+2} - F_{k+1}$ lalu amati jumlahnya runtuh. Periksalah
        untuk $n = 6$.
  \item Buktikan identitas kuadratnya
        $F_1^2 + F_2^2 + \dots + F_n^2 = F_n F_{n+1}$,
        secara teleskopik dengan
        $F_k F_{k+1} - F_{k-1} F_k = F_k^2$. Periksalah untuk
        $n = 4$. (Gambarannya: persegi bersisi
        $1, 1, 2, 3, 5, \dots$ memubin sebuah persegi panjang --- kerangka
        spiral Fibonacci yang termasyhur.)
  \item Identitas Cassini menyatakan
        $F_{n+1} F_{n-1} - F_n^2 = (-1)^n$. Periksalah untuk
        $n = 4, 5, 6$ --- lalu kenali mesin di balik
        kiat persegi yang lenyap pada soal luas di jilid
        sebelumnya.
  \item Tunjukkan dari rekurensinya bahwa
        $F_{n+2} \geq 2 F_n$: Fibonacci paling sedikit berlipat dua setiap
        dua langkah --- jadi tumbuh paling sedikit secepat \omterm{def:g11:seq:geometric}{barisan geometri}
        berasio $\sqrt2$.
  \item Hitunglah \omterm{def:g11:seq:geometric}{rasio} $r_n = \frac{F_{n+1}}{F_n}$ untuk
        $n = 3$ sampai $10$ (tiga angka desimal). Dengan menerima bahwa
        \omterm{def:g11:seq:geometric}{rasio} itu mengendap pada limit $L$, terapkan hubungan
        $r_{n+1} = 1 + \frac{1}{r_n}$ pada limitnya lalu selesaikan:
        bilangan mana pada \cref{pb:g10:algebra:1} yang dipuja para
        kelinci itu?
\end{enumerate}

\medskip
\noindent\textbf{Bagian III --- Kiat \omterm{def:g11:seq:sequence}{barisan} bantu, di
bank.}
\begin{enumerate}[resume]
  \item Perumumlah \cref{exo:g11:seq:11}: untuk
        $u_{n+1} = a\,u_n + b$ dengan $a \neq 1$, ambil
        $\ell = \frac{b}{1 - a}$ (\omterm{pb:g10:functions:1}{titik tetapnya}). Tunjukkan bahwa
        $v_n = u_n - \ell$ \omterm{def:g11:seq:geometric}{geometri} dengan \omterm{def:g11:seq:geometric}{rasio} $a$, lalu
        simpulkan
        $u_n = a^n (u_0 - \ell) + \ell$.
  \item Sebuah pinjaman: $10\,000$ euro dengan bunga $1\,\%$ per bulan,
        diangsur $300$ euro per bulan, sehingga utangnya menuruti
        $d_{n+1} = 1.01\,d_n - 300$. Terapkan pertanyaan 13 (\omterm{pb:g10:functions:1}{titik
        tetapnya} dahulu!) untuk memperoleh rumus eksplisit $d_n$.
  \item Dengan kalkulator, carilah bulan pertama saat
        utangnya lunas, dan jumlah seluruh angsurannya. Berapa
        ongkos meminjam itu sendiri?
  \item Sebuah kota berpenduduk $50\,000$ jiwa tumbuh $2\,\%$ setahun
        dan menerima $1\,000$ pendatang baru selain itu:
        $p_{n+1} = 1.02\,p_n + 1000$. Berikan rumus eksplisitnya
        dan penduduknya setelah $10$ tahun.
\end{enumerate}

\medskip
\noindent\textbf{Bagian IV --- Dua keluarga bangsawan.}
\begin{enumerate}[resume]
  \item Hitunglah $1 + 2 + 3 + \dots + 1000$
        (\cref{thm:g11:seq:intsum} --- jumlah si kecil Gauss dari
        jilid sebelumnya, kini resmi), lalu
        $1 + 2 + 4 + \dots + 2^{19}$
        (\cref{thm:g11:seq:geomsum}).
  \item Hitunglah jumlah \omterm{def:g11:seq:arithmetic}{barisan aritmetika}
        $7, 12, 17, \dots, 502$ (berapa sukunya?).
  \item Rencana tabungan: $100$ euro disetorkan setiap bulan, memperoleh
        $0.5\,\%$ per bulan; setelah setoran ke-$n$ saldonya
        adalah $100\left(1.005^{n-1} + \dots + 1.005 + 1\right)$.
        Hitunglah saldonya setelah $5$ tahun ($n = 60$).
  \item Penutup --- kotak perkakas penjinak \omterm{def:g11:seq:sequence}{barisan}: perian eksplisit
        melawan perian rekuren; dua keluarga bangsawan dan
        rumus jumlahnya; \omterm{def:g11:seq:sequence}{barisan} bantu yang mengubah
        rekurensi afin menjadi rekurensi \omterm{def:g11:seq:geometric}{geometri}; serta Fibonacci,
        warga pertama di luar kedua keluarga itu, yang hari ini dijinakkan
        oleh identitas sambil menanti matriks (tahun berikutnya) dan limit
        untuk ditangkap sepenuhnya. Masing-masing satu kalimat.

\end{enumerate}
\end{problem}
